% !TeX root = ../Western Philosophy.tex
\section{帕斯卡和马勒布朗士}

\subsection{Pascal哲学}

\subsubsection{生平}

1623年生，青年科学英才，16岁发表圆锥曲线论，发明物理机械计算器。在气压方面和真空方面有卓越贡献。
（注释：Descart否认原子论的虚空，同时否认空间连续情况下没有物质的真空）

Pascal本人的研究中心是神学，1657年作《致外省人信札》。Pascal有詹森主义者的影响，展现出詹森主义的思想气质，1662年去世，1670年出版《思想录》。

\subsubsection{名言}
哲思

\textbf{The external silence od the infinite spaces terrifies me}

\textbf{We die alone}

历史

\textbf{Had Cleopatra's nose been shorter, the whole face of the world would have been changed}

\textbf{The heart has its reasons of which reason knows nothing} （直观的决定并不完全依赖于理性（后一个reason是官能，前一个reason是行动；heart是理智上直观的东西））

大名句：\textbf{Man is only a reed, the frailest thing in the world, but he is a thinking reed}

\subsubsection{哲学观点}

\paragraph{批评耶稣会的教义}

坚持岗哨的自由意志（参考：兼容论vs自由意志论）与Hobbes的机械性论证不同，Pascal认为这些will来自于上帝的恩典（詹森主义者认为的恩典）。

坚持严格的道德准则，反对含糊其词，可能主义（可以选择较为宽松的那个要求），意愿方向等“套路”。（耶稣会认为心中想的东西是诚实的嘴上骗人不算撒谎。。。？）

\paragraph{Pascal’s Wager帕斯卡之赌}

理性无法证明上帝存在（要靠heart直接了解和看到）。Pascal的理性更多指推论，和Descart并不完全冲突。稍微弱一点的情况下，可以用推论证明上帝的存在，于是提出了以下赌局：

\begin{center}
\begin{tabular}{lcc}
\toprule
 & 上帝存在 & 上帝不存在 \\
\midrule
信 & 幸福$\infty$ & 有限损失（like-10000） \\
不信 & 痛苦$\infty$ & 有限收益 （like+10000）\\
\bottomrule
\end{tabular}
\end{center}

显然有很大的问题，比如信错了教的危害远远大于不信任何教）（我为什么不能信菩萨呢）
